\documentclass[11pt,a4paper]{article}
\usepackage[utf8]{inputenc}
\usepackage{amsmath,amssymb}
\usepackage{graphicx}
\usepackage{booktabs}
\usepackage{hyperref}

\title{Cinemática Directa e Inversa de Manipuladores Industriales de 6 Grados de Libertad (6-DOF)}
\author{Dr. Carlos E. Morales \and Ing. Mariana Gómez \\ 
\textit{Laboratorio de Robótica y Automatización Industrial, Centro Nova}}
\date{Septiembre de 2026}

\begin{document}

\maketitle

\begin{abstract}
Este documento técnico presenta la formulación cinemática rigurosa para brazos robóticos industriales de 6 grados de libertad (6-DOF) con muñeca esférica ortogonal, tales como los manipuladores ABB IRB 1600, FANUC M-20iA y KUKA KR Quantec. Se establece la convención clásica de Denavit-Hartenberg (D-H), deduciendo las matrices de transformación homogénea entre eslabones consecutivos. Posteriormente, se resuelve el problema de cinemática inversa mediante el método de desacoplo cinemático de Pieper, separando la posición del centro de la muñeca de su orientación angular. Finalmente, se deriva la matriz Jacobiana analítica y se determinan las configuraciones singulares críticas para evitar inestabilidad en actuadores servocontrolados.
\end{abstract}

\section{Introducción y Arquitectura del Manipulador}
Los manipuladores robóticos industriales de seis ejes articulados constituyen el estándar fundamental en líneas de manufactura automotriz, soldadura MIG/TIG y manipulación de precisión. La estructura cinemática típica se compone de:
\begin{itemize}
    \item \textbf{Cintura (Joint 1):} Rotación en el plano horizontal respecto a la base fija ($z_0$).
    \item \textbf{Hombro (Joint 2):} Rotación en el plano sagital que eleva el brazo superior.
    \item \textbf{Codo (Joint 3):} Articulación planar paralela al eje del hombro.
    \item \textbf{Muñeca esférica (Joints 4, 5 y 6):} Tres ejes de rotación concurrentes en un único punto focal $P_w$.
\end{itemize}

El espacio de estados de las articulaciones se define por el vector de coordenadas generalizadas:
\begin{equation}
q = [\theta_1, \theta_2, \theta_3, \theta_4, \theta_5, \theta_6]^T \in \mathbb{R}^6
\end{equation}

\section{Convención de Parámetros Denavit-Hartenberg (D-H)}
Para establecer el sistema de referencia solidario a cada eslabón $i \in \{1,\dots,6\}$, se aplican las reglas geométricas de Denavit-Hartenberg basadas en cuatro parámetros:
\begin{enumerate}
    \item $\theta_i$: Ángulo de rotación alrededor del eje $z_{i-1}$.
    \item $d_i$: Desplazamiento lineal a lo largo del eje $z_{i-1}$.
    \item $a_i$: Longitud del eslabón (distancia a lo largo de $x_i$).
    \item $\alpha_i$: Ángulo de torsión entre $z_{i-1}$ y $z_i$ respecto a $x_i$.
\end{enumerate}

\subsection{Tabla de Parámetros Geométricos}
A continuación se detalla la tabla D-H estándar para un robot de 6-DOF con longitudes de eslabón $d_1 = 450\,\text{mm}$, $a_2 = 400\,\text{mm}$, $a_3 = 135\,\text{mm}$, $d_4 = 420\,\text{mm}$ y $d_6 = 100\,\text{mm}$:

\begin{table}[h]
\centering
\begin{tabular}{ccccc}
\hline
\textbf{Articulación $i$} & $\theta_i$ [rad] & $d_i$ [mm] & $a_i$ [mm] & $\alpha_i$ [rad] \\
\hline
1 & $\theta_1^*$ & $d_1$ & 0 & $\pi/2$ \\
2 & $\theta_2^*$ & 0 & $a_2$ & 0 \\
3 & $\theta_3^*$ & 0 & $a_3$ & $\pi/2$ \\
4 & $\theta_4^*$ & $d_4$ & 0 & $-\pi/2$ \\
5 & $\theta_5^*$ & 0 & 0 & $\pi/2$ \\
6 & $\theta_6^*$ & $d_6$ & 0 & 0 \\
\hline
\end{tabular}
\end{table}

\section{Cinemática Directa y Matrices de Transformación}
La matriz de transformación homogénea del sistema $\{i\}$ respecto a $\{i-1\}$ se formula mediante la multiplicación de rotaciones y traslaciones básicas:
\begin{equation}
A_i = \text{Rot}(z, \theta_i) \cdot \text{Trans}(z, d_i) \cdot \text{Trans}(x, a_i) \cdot \text{Rot}(x, \alpha_i)
\end{equation}

Desarrollando los productos trigonométricos compactos con $c_i = \cos(\theta_i)$ y $s_i = \sin(\theta_i)$:
\begin{equation}
A_i = \begin{bmatrix}
\cos\theta_i & -\sin\theta_i \cos\alpha_i & \sin\theta_i \sin\alpha_i & a_i \cos\theta_i \\
\sin\theta_i & \cos\theta_i \cos\alpha_i & -\cos\theta_i \sin\alpha_i & a_i \sin\theta_i \\
0 & \sin\alpha_i & \cos\alpha_i & d_i \\
0 & 0 & 0 & 1
\end{bmatrix}
\end{equation}

La postura global del efector final respecto a la base inercial $T_6^0$ se calcula como el producto encadenado:
\begin{equation}
T_6^0(q) = A_1(\theta_1) \cdot A_2(\theta_2) \cdot A_3(\theta_3) \cdot A_4(\theta_4) \cdot A_5(\theta_5) \cdot A_6(\theta_6) = \begin{bmatrix}
R_{3\times3} & p_{3\times1} \\
0_{1\times3} & 1
\end{bmatrix}
\end{equation}
donde $R$ describe la orientación tridimensional del extremo del robot y $p = [p_x, p_y, p_z]^T$ denota la posición cartesiana.

\section{Cinemática Inversa: Método de Pieper}
Dado el efector final en una posición de referencia deseada $(p, R)$, el cálculo de los ángulos articulares $q$ requiere resolver un sistema no lineal de ecuaciones trigonométricas acopladas. Gracias a la intersección concurrente de los últimos tres ejes (muñeca esférica), el problema se desacopla en dos fases analíticas independientes.

\subsection{Paso 1: Posición del Centro de la Muñeca}
El centro focal de la muñeca $p_w$ se aísla restando la traslación de la herramienta $d_6$ a lo largo del vector director $z_6$:
\begin{equation}
p_w = p - d_6 \cdot R \cdot \begin{bmatrix} 0 \\ 0 \\ 1 \end{bmatrix} = \begin{bmatrix} x_w \\ y_w \\ z_w \end{bmatrix}
\end{equation}

El ángulo de la cintura $\theta_1$ se determina directamente en el plano cartesiano base:
\begin{equation}
\theta_1 = \text{atan2}(y_w, x_w) \quad \text{o} \quad \theta_1 = \text{atan2}(y_w, x_w) + \pi
\end{equation}

Posteriormente, proyectando en el plano del brazo, se calculan $\theta_2$ y $\theta_3$ aplicando el teorema del coseno:
\begin{equation}
\cos\theta_3 = \frac{(x_w c_1 + y_w s_1)^2 + (z_w - d_1)^2 - a_2^2 - (a_3^2 + d_4^2)}{2 a_2 \sqrt{a_3^2 + d_4^2}}
\end{equation}

\subsection{Paso 2: Orientación de la Muñeca Esférica}
Una vez fijados $(\theta_1, \theta_2, \theta_3)$, se computa la orientación relativa que debe proveer la muñeca:
\begin{equation}
R_{3\times3}^{wrist} = (R_3^0)^T \cdot R_{target} = R_6^3(\theta_4, \theta_5, \theta_6)
\end{equation}

Este producto matricial equivale a una secuencia de ángulos de Euler $Z-Y-Z$, de donde se derivan:
\begin{equation}
\theta_5 = \text{atan2}\left(\pm\sqrt{(R_{13}^{wrist})^2 + (R_{23}^{wrist})^2}, R_{33}^{wrist}\right)
\end{equation}

\section{Matriz Jacobiana y Singularidades Cinemáticas}
La relación entre las velocidades articulares $\dot{q}$ y las velocidades espaciales del efector final $v = [\dot{p}^T, \omega^T]^T$ se rige por la matriz Jacobiana geométrica:
\begin{equation}
v = J(q) \cdot \dot{q} = \begin{bmatrix} J_v(q) \\ J_\omega(q) \end{bmatrix} \dot{q}
\end{equation}

Las singularidades cinemáticas ocurren cuando $\det(J(q)) = 0$, provocando pérdida instantánea de uno o más grados de libertad y demandando velocidades articulares infinitas para seguir trayectorias cartesianas:
\begin{itemize}
    \item \textbf{Singularidad de Muñeca ($\theta_5 = 0$):} Los ejes $z_4$ y $z_6$ quedan alineados, colapsando un grado de orientación.
    \item \textbf{Singularidad de Frontera:} El manipulador alcanza su extensión máxima radial, saliendo del volumen de trabajo útil.
    \item \textbf{Singularidad de Codo:} El centro de la muñeca $p_w$ intersecta el eje $z_1$, eliminando la capacidad de movimiento tangencial.
\end{itemize}

\section{Conclusiones}
La formulación analítica de la cinemática directa e inversa mediante el criterio de Denavit-Hartenberg y el desacoplo de Pieper permite una ejecución en tiempo real en controladores PLC industriales con ciclos de reloj inferiores a $1\,\text{ms}$, garantizando alta repetibilidad en manufactura industrial.

\begin{thebibliography}{9}
\bibitem{craig} J. J. Craig, \textit{Introduction to Robotics: Mechanics and Control}, 4th ed., Pearson, 2018.
\bibitem{spong} M. W. Spong, S. Hutchinson, and M. Vidyasagar, \textit{Robot Modeling and Control}, John Wiley \& Sons, 2020.
\bibitem{siciliano} B. Siciliano and O. Khatib, \textit{Springer Handbook of Robotics}, 2nd ed., Springer, 2016.
\end{thebibliography}

\end{document}
